\section{Introduction}

<field context -> concrete problem -> technical constraint; each factual sentence with ~\cite{key} or >

\section{Related Work}

\subsection{Large language models}

<compare these works on the group's axis: findings, conditions, limits ~\cite{k1,k2}>

\subsection{Modelling the scaling behavior}

<compare these works on the group's axis: findings, conditions, limits ~\cite{k1,k2}>

\subsection{Estimating hyperparameters for large models}

<compare these works on the group's axis: findings, conditions, limits ~\cite{k1,k2}>

\subsection{Improved model architectures}

<compare these works on the group's axis: findings, conditions, limits ~\cite{k1,k2}>

\subsection{Efficient frontier}

<compare these works on the group's axis: findings, conditions, limits ~\cite{k1,k2}>

<gap derived from the shared limitation of the works above ~\cite{k}>

Our contributions are:

\begin{itemize}
\item We find that current large language models are significantly undertrained, a consequence of the recent focus on scaling language models whilst keeping the amount of training data constant.
\item By training over 400 language models ranging from 70 million to over 16 billion parameters on 5 to 500 billion tokens, we find that for compute-optimal training, the model size and the number of training tokens should be scaled equally: for every doubling of model size the number of training tokens …
\item Instead, we find that model size and the number of training tokens should be scaled in equal proportions.
\item We find that all three methods predict that current large models should be substantially smaller and therefore trained much longer than is currently done.
\end{itemize}
